\section{未锚定账本事件的叙事补充：mongol-yuan}

\label{sec:anchor-batch-two-mongol-yuan}

本组为尚无正文跳转目标的账本事件补入锚点。每一记录置于战争、行政、环境、迁徙与地方社会的关系中，并按来源限度约束叙事。

\subsection{制度、区域与材料边界}

\eventanchor{YUAN-1269-REGNAL-YEAR}{1269年·蒙古／元以忽必烈在位第10年作为诏令、赋役与外交文书的纪年框架}账本条目记录蒙古／元在1269年的“蒙古／元以忽必烈在位第10年作为诏令、赋役与外交文书的纪年框架”。其背景是新岁颁历、年号和纪年是中枢与地方行政沟通的共同前提，后续影响为为同年政令、战事和地方材料提供日期锚点，不能单独推知具体政策执行。政权、军队或制度的变化须经由道路、文书、资源和地方中介才会落实，城市、乡村、边地和流动人群不可能在同一时间承受相同结果。

本条使用《元史》〈本纪〉及相关志；《元朝秘史》；《元典章》或《通制条格》相关条，并标示“纪年框架可靠；该记录仅确证政权纪年连续，年内具体事项须由本年条另证”。这些材料能够钉住年序、命令、战役或局部地点，却不能直接推出人口、税收、身份和社会动机的总量。应区分诏令与实施、条约与边境生活、军报与伤亡统计，也不将官府的族群或行政称谓写成固定的社会本质。

从区域史角度看，该事件是资源、信息与权力重新连接的节点。不同地方会因市场、交通、宗教组织、家庭策略和环境条件而产生不同反应；材料的沉默限制我们对未被记录者所能作出的判断。

\eventanchor{YUAN-1270-REGNAL-YEAR}{1270年·蒙古／元以忽必烈在位第11年作为诏令、赋役与外交文书的纪年框架}账本条目记录蒙古／元在1270年的“蒙古／元以忽必烈在位第11年作为诏令、赋役与外交文书的纪年框架”。其背景是新岁颁历、年号和纪年是中枢与地方行政沟通的共同前提，后续影响为为同年政令、战事和地方材料提供日期锚点，不能单独推知具体政策执行。政权、军队或制度的变化须经由道路、文书、资源和地方中介才会落实，城市、乡村、边地和流动人群不可能在同一时间承受相同结果。

本条使用《元史》〈本纪〉及相关志；《元朝秘史》；《元典章》或《通制条格》相关条，并标示“纪年框架可靠；该记录仅确证政权纪年连续，年内具体事项须由本年条另证”。这些材料能够钉住年序、命令、战役或局部地点，却不能直接推出人口、税收、身份和社会动机的总量。应区分诏令与实施、条约与边境生活、军报与伤亡统计，也不将官府的族群或行政称谓写成固定的社会本质。

从区域史角度看，该事件是资源、信息与权力重新连接的节点。不同地方会因市场、交通、宗教组织、家庭策略和环境条件而产生不同反应；材料的沉默限制我们对未被记录者所能作出的判断。

\eventanchor{YUAN-1271-EVENT-173}{1271年·忽必烈定国号元}账本条目记录元在1271年的“忽必烈定国号元”。其背景是前期边防、财政和统治联盟的压力构成直接背景，具体条件须按本条材料分辨，后续影响为改变后续资源配置与政治选择；实际执行范围和社会影响仍须分项核验。政权、军队或制度的变化须经由道路、文书、资源和地方中介才会落实，城市、乡村、边地和流动人群不可能在同一时间承受相同结果。

本条使用《元史》〈本纪〉及相关志；《元朝秘史》；《元典章》或《通制条格》相关条，并标示“编年主干可靠；细节、规模与动机须与对方材料、地方文书或考古材料互校”。这些材料能够钉住年序、命令、战役或局部地点，却不能直接推出人口、税收、身份和社会动机的总量。应区分诏令与实施、条约与边境生活、军报与伤亡统计，也不将官府的族群或行政称谓写成固定的社会本质。

从区域史角度看，该事件是资源、信息与权力重新连接的节点。不同地方会因市场、交通、宗教组织、家庭策略和环境条件而产生不同反应；材料的沉默限制我们对未被记录者所能作出的判断。

\eventanchor{YUAN-1271-REGNAL-YEAR}{1271年·蒙古／元以忽必烈在位第12年作为诏令、赋役与外交文书的纪年框架}账本条目记录蒙古／元在1271年的“蒙古／元以忽必烈在位第12年作为诏令、赋役与外交文书的纪年框架”。其背景是新岁颁历、年号和纪年是中枢与地方行政沟通的共同前提，后续影响为为同年政令、战事和地方材料提供日期锚点，不能单独推知具体政策执行。政权、军队或制度的变化须经由道路、文书、资源和地方中介才会落实，城市、乡村、边地和流动人群不可能在同一时间承受相同结果。

本条使用《元史》〈本纪〉及相关志；《元朝秘史》；《元典章》或《通制条格》相关条，并标示“纪年框架可靠；该记录仅确证政权纪年连续，年内具体事项须由本年条另证”。这些材料能够钉住年序、命令、战役或局部地点，却不能直接推出人口、税收、身份和社会动机的总量。应区分诏令与实施、条约与边境生活、军报与伤亡统计，也不将官府的族群或行政称谓写成固定的社会本质。

从区域史角度看，该事件是资源、信息与权力重新连接的节点。不同地方会因市场、交通、宗教组织、家庭策略和环境条件而产生不同反应；材料的沉默限制我们对未被记录者所能作出的判断。

\eventanchor{YUAN-1272-REGNAL-YEAR}{1272年·蒙古／元以忽必烈在位第13年作为诏令、赋役与外交文书的纪年框架}账本条目记录蒙古／元在1272年的“蒙古／元以忽必烈在位第13年作为诏令、赋役与外交文书的纪年框架”。其背景是新岁颁历、年号和纪年是中枢与地方行政沟通的共同前提，后续影响为为同年政令、战事和地方材料提供日期锚点，不能单独推知具体政策执行。政权、军队或制度的变化须经由道路、文书、资源和地方中介才会落实，城市、乡村、边地和流动人群不可能在同一时间承受相同结果。

本条使用《元史》〈本纪〉及相关志；《元朝秘史》；《元典章》或《通制条格》相关条，并标示“纪年框架可靠；该记录仅确证政权纪年连续，年内具体事项须由本年条另证”。这些材料能够钉住年序、命令、战役或局部地点，却不能直接推出人口、税收、身份和社会动机的总量。应区分诏令与实施、条约与边境生活、军报与伤亡统计，也不将官府的族群或行政称谓写成固定的社会本质。

从区域史角度看，该事件是资源、信息与权力重新连接的节点。不同地方会因市场、交通、宗教组织、家庭策略和环境条件而产生不同反应；材料的沉默限制我们对未被记录者所能作出的判断。

\eventanchor{YUAN-1273-EVENT-174}{1273年·襄阳、樊城陷落}账本条目记录元与南宋在1273年的“襄阳、樊城陷落”。其背景是前期边防、财政和统治联盟的压力构成直接背景，具体条件须按本条材料分辨，后续影响为改变后续资源配置与政治选择；实际执行范围和社会影响仍须分项核验。政权、军队或制度的变化须经由道路、文书、资源和地方中介才会落实，城市、乡村、边地和流动人群不可能在同一时间承受相同结果。

本条使用《元史》〈本纪〉及相关志；《元朝秘史》；《元典章》或《通制条格》相关条，并标示“编年主干可靠；细节、规模与动机须与对方材料、地方文书或考古材料互校”。这些材料能够钉住年序、命令、战役或局部地点，却不能直接推出人口、税收、身份和社会动机的总量。应区分诏令与实施、条约与边境生活、军报与伤亡统计，也不将官府的族群或行政称谓写成固定的社会本质。

从区域史角度看，该事件是资源、信息与权力重新连接的节点。不同地方会因市场、交通、宗教组织、家庭策略和环境条件而产生不同反应；材料的沉默限制我们对未被记录者所能作出的判断。

\eventanchor{YUAN-1273-REGNAL-YEAR}{1273年·蒙古／元以忽必烈在位第14年作为诏令、赋役与外交文书的纪年框架}账本条目记录蒙古／元在1273年的“蒙古／元以忽必烈在位第14年作为诏令、赋役与外交文书的纪年框架”。其背景是新岁颁历、年号和纪年是中枢与地方行政沟通的共同前提，后续影响为为同年政令、战事和地方材料提供日期锚点，不能单独推知具体政策执行。政权、军队或制度的变化须经由道路、文书、资源和地方中介才会落实，城市、乡村、边地和流动人群不可能在同一时间承受相同结果。

本条使用《元史》〈本纪〉及相关志；《元朝秘史》；《元典章》或《通制条格》相关条，并标示“纪年框架可靠；该记录仅确证政权纪年连续，年内具体事项须由本年条另证”。这些材料能够钉住年序、命令、战役或局部地点，却不能直接推出人口、税收、身份和社会动机的总量。应区分诏令与实施、条约与边境生活、军报与伤亡统计，也不将官府的族群或行政称谓写成固定的社会本质。

从区域史角度看，该事件是资源、信息与权力重新连接的节点。不同地方会因市场、交通、宗教组织、家庭策略和环境条件而产生不同反应；材料的沉默限制我们对未被记录者所能作出的判断。

\eventanchor{YUAN-1274-EVENT-175}{1274年·元军分路南下进攻南宋}账本条目记录元与南宋在1274年的“元军分路南下进攻南宋”。其背景是前期边防、财政和统治联盟的压力构成直接背景，具体条件须按本条材料分辨，后续影响为改变后续资源配置与政治选择；实际执行范围和社会影响仍须分项核验。政权、军队或制度的变化须经由道路、文书、资源和地方中介才会落实，城市、乡村、边地和流动人群不可能在同一时间承受相同结果。

本条使用《元史》〈本纪〉及相关志；《元朝秘史》；《元典章》或《通制条格》相关条，并标示“编年主干可靠；细节、规模与动机须与对方材料、地方文书或考古材料互校”。这些材料能够钉住年序、命令、战役或局部地点，却不能直接推出人口、税收、身份和社会动机的总量。应区分诏令与实施、条约与边境生活、军报与伤亡统计，也不将官府的族群或行政称谓写成固定的社会本质。

从区域史角度看，该事件是资源、信息与权力重新连接的节点。不同地方会因市场、交通、宗教组织、家庭策略和环境条件而产生不同反应；材料的沉默限制我们对未被记录者所能作出的判断。

\eventanchor{YUAN-1274-REGNAL-YEAR}{1274年·蒙古／元以忽必烈在位第15年作为诏令、赋役与外交文书的纪年框架}账本条目记录蒙古／元在1274年的“蒙古／元以忽必烈在位第15年作为诏令、赋役与外交文书的纪年框架”。其背景是新岁颁历、年号和纪年是中枢与地方行政沟通的共同前提，后续影响为为同年政令、战事和地方材料提供日期锚点，不能单独推知具体政策执行。政权、军队或制度的变化须经由道路、文书、资源和地方中介才会落实，城市、乡村、边地和流动人群不可能在同一时间承受相同结果。

本条使用《元史》〈本纪〉及相关志；《元朝秘史》；《元典章》或《通制条格》相关条，并标示“纪年框架可靠；该记录仅确证政权纪年连续，年内具体事项须由本年条另证”。这些材料能够钉住年序、命令、战役或局部地点，却不能直接推出人口、税收、身份和社会动机的总量。应区分诏令与实施、条约与边境生活、军报与伤亡统计，也不将官府的族群或行政称谓写成固定的社会本质。

从区域史角度看，该事件是资源、信息与权力重新连接的节点。不同地方会因市场、交通、宗教组织、家庭策略和环境条件而产生不同反应；材料的沉默限制我们对未被记录者所能作出的判断。

\eventanchor{YUAN-1275-REGNAL-YEAR}{1275年·蒙古／元以忽必烈在位第16年作为诏令、赋役与外交文书的纪年框架}账本条目记录蒙古／元在1275年的“蒙古／元以忽必烈在位第16年作为诏令、赋役与外交文书的纪年框架”。其背景是新岁颁历、年号和纪年是中枢与地方行政沟通的共同前提，后续影响为为同年政令、战事和地方材料提供日期锚点，不能单独推知具体政策执行。政权、军队或制度的变化须经由道路、文书、资源和地方中介才会落实，城市、乡村、边地和流动人群不可能在同一时间承受相同结果。

本条使用《元史》〈本纪〉及相关志；《元朝秘史》；《元典章》或《通制条格》相关条，并标示“纪年框架可靠；该记录仅确证政权纪年连续，年内具体事项须由本年条另证”。这些材料能够钉住年序、命令、战役或局部地点，却不能直接推出人口、税收、身份和社会动机的总量。应区分诏令与实施、条约与边境生活、军报与伤亡统计，也不将官府的族群或行政称谓写成固定的社会本质。

从区域史角度看，该事件是资源、信息与权力重新连接的节点。不同地方会因市场、交通、宗教组织、家庭策略和环境条件而产生不同反应；材料的沉默限制我们对未被记录者所能作出的判断。

\eventanchor{YUAN-1276-EVENT-176}{1276年·元军进入临安，宋恭帝降元}账本条目记录元与南宋在1276年的“元军进入临安，宋恭帝降元”。其背景是前期边防、财政和统治联盟的压力构成直接背景，具体条件须按本条材料分辨，后续影响为改变后续资源配置与政治选择；实际执行范围和社会影响仍须分项核验。政权、军队或制度的变化须经由道路、文书、资源和地方中介才会落实，城市、乡村、边地和流动人群不可能在同一时间承受相同结果。

本条使用《元史》〈本纪〉及相关志；《元朝秘史》；《元典章》或《通制条格》相关条，并标示“编年主干可靠；细节、规模与动机须与对方材料、地方文书或考古材料互校”。这些材料能够钉住年序、命令、战役或局部地点，却不能直接推出人口、税收、身份和社会动机的总量。应区分诏令与实施、条约与边境生活、军报与伤亡统计，也不将官府的族群或行政称谓写成固定的社会本质。

从区域史角度看，该事件是资源、信息与权力重新连接的节点。不同地方会因市场、交通、宗教组织、家庭策略和环境条件而产生不同反应；材料的沉默限制我们对未被记录者所能作出的判断。

\eventanchor{YUAN-1276-REGNAL-YEAR}{1276年·蒙古／元以忽必烈在位第17年作为诏令、赋役与外交文书的纪年框架}账本条目记录蒙古／元在1276年的“蒙古／元以忽必烈在位第17年作为诏令、赋役与外交文书的纪年框架”。其背景是新岁颁历、年号和纪年是中枢与地方行政沟通的共同前提，后续影响为为同年政令、战事和地方材料提供日期锚点，不能单独推知具体政策执行。政权、军队或制度的变化须经由道路、文书、资源和地方中介才会落实，城市、乡村、边地和流动人群不可能在同一时间承受相同结果。

本条使用《元史》〈本纪〉及相关志；《元朝秘史》；《元典章》或《通制条格》相关条，并标示“纪年框架可靠；该记录仅确证政权纪年连续，年内具体事项须由本年条另证”。这些材料能够钉住年序、命令、战役或局部地点，却不能直接推出人口、税收、身份和社会动机的总量。应区分诏令与实施、条约与边境生活、军报与伤亡统计，也不将官府的族群或行政称谓写成固定的社会本质。

从区域史角度看，该事件是资源、信息与权力重新连接的节点。不同地方会因市场、交通、宗教组织、家庭策略和环境条件而产生不同反应；材料的沉默限制我们对未被记录者所能作出的判断。

\eventanchor{YUAN-1277-EVENT-177}{1277年·元在江南建立行省、路府与财赋接收安排}账本条目记录元在1277年的“元在江南建立行省、路府与财赋接收安排”。其背景是前期边防、财政和统治联盟的压力构成直接背景，具体条件须按本条材料分辨，后续影响为改变后续资源配置与政治选择；实际执行范围和社会影响仍须分项核验。政权、军队或制度的变化须经由道路、文书、资源和地方中介才会落实，城市、乡村、边地和流动人群不可能在同一时间承受相同结果。

本条使用《元史》〈本纪〉及相关志；《元朝秘史》；《元典章》或《通制条格》相关条，并标示“编年主干可靠；细节、规模与动机须与对方材料、地方文书或考古材料互校”。这些材料能够钉住年序、命令、战役或局部地点，却不能直接推出人口、税收、身份和社会动机的总量。应区分诏令与实施、条约与边境生活、军报与伤亡统计，也不将官府的族群或行政称谓写成固定的社会本质。

从区域史角度看，该事件是资源、信息与权力重新连接的节点。不同地方会因市场、交通、宗教组织、家庭策略和环境条件而产生不同反应；材料的沉默限制我们对未被记录者所能作出的判断。

\eventanchor{YUAN-1277-REGNAL-YEAR}{1277年·蒙古／元以忽必烈在位第18年作为诏令、赋役与外交文书的纪年框架}账本条目记录蒙古／元在1277年的“蒙古／元以忽必烈在位第18年作为诏令、赋役与外交文书的纪年框架”。其背景是新岁颁历、年号和纪年是中枢与地方行政沟通的共同前提，后续影响为为同年政令、战事和地方材料提供日期锚点，不能单独推知具体政策执行。政权、军队或制度的变化须经由道路、文书、资源和地方中介才会落实，城市、乡村、边地和流动人群不可能在同一时间承受相同结果。

本条使用《元史》〈本纪〉及相关志；《元朝秘史》；《元典章》或《通制条格》相关条，并标示“纪年框架可靠；该记录仅确证政权纪年连续，年内具体事项须由本年条另证”。这些材料能够钉住年序、命令、战役或局部地点，却不能直接推出人口、税收、身份和社会动机的总量。应区分诏令与实施、条约与边境生活、军报与伤亡统计，也不将官府的族群或行政称谓写成固定的社会本质。

从区域史角度看，该事件是资源、信息与权力重新连接的节点。不同地方会因市场、交通、宗教组织、家庭策略和环境条件而产生不同反应；材料的沉默限制我们对未被记录者所能作出的判断。

\eventanchor{YUAN-1278-REGNAL-YEAR}{1278年·蒙古／元以忽必烈在位第19年作为诏令、赋役与外交文书的纪年框架}账本条目记录蒙古／元在1278年的“蒙古／元以忽必烈在位第19年作为诏令、赋役与外交文书的纪年框架”。其背景是新岁颁历、年号和纪年是中枢与地方行政沟通的共同前提，后续影响为为同年政令、战事和地方材料提供日期锚点，不能单独推知具体政策执行。政权、军队或制度的变化须经由道路、文书、资源和地方中介才会落实，城市、乡村、边地和流动人群不可能在同一时间承受相同结果。

本条使用《元史》〈本纪〉及相关志；《元朝秘史》；《元典章》或《通制条格》相关条，并标示“纪年框架可靠；该记录仅确证政权纪年连续，年内具体事项须由本年条另证”。这些材料能够钉住年序、命令、战役或局部地点，却不能直接推出人口、税收、身份和社会动机的总量。应区分诏令与实施、条约与边境生活、军报与伤亡统计，也不将官府的族群或行政称谓写成固定的社会本质。

从区域史角度看，该事件是资源、信息与权力重新连接的节点。不同地方会因市场、交通、宗教组织、家庭策略和环境条件而产生不同反应；材料的沉默限制我们对未被记录者所能作出的判断。

\eventanchor{YUAN-1279-EVENT-178}{1279年·崖山海战后南宋政权终结}账本条目记录元与南宋在1279年的“崖山海战后南宋政权终结”。其背景是前期边防、财政和统治联盟的压力构成直接背景，具体条件须按本条材料分辨，后续影响为改变后续资源配置与政治选择；实际执行范围和社会影响仍须分项核验。政权、军队或制度的变化须经由道路、文书、资源和地方中介才会落实，城市、乡村、边地和流动人群不可能在同一时间承受相同结果。

本条使用《元史》〈本纪〉及相关志；《元朝秘史》；《元典章》或《通制条格》相关条，并标示“编年主干可靠；细节、规模与动机须与对方材料、地方文书或考古材料互校”。这些材料能够钉住年序、命令、战役或局部地点，却不能直接推出人口、税收、身份和社会动机的总量。应区分诏令与实施、条约与边境生活、军报与伤亡统计，也不将官府的族群或行政称谓写成固定的社会本质。

从区域史角度看，该事件是资源、信息与权力重新连接的节点。不同地方会因市场、交通、宗教组织、家庭策略和环境条件而产生不同反应；材料的沉默限制我们对未被记录者所能作出的判断。

\eventanchor{YUAN-1279-REGNAL-YEAR}{1279年·蒙古／元以忽必烈在位第20年作为诏令、赋役与外交文书的纪年框架}账本条目记录蒙古／元在1279年的“蒙古／元以忽必烈在位第20年作为诏令、赋役与外交文书的纪年框架”。其背景是新岁颁历、年号和纪年是中枢与地方行政沟通的共同前提，后续影响为为同年政令、战事和地方材料提供日期锚点，不能单独推知具体政策执行。政权、军队或制度的变化须经由道路、文书、资源和地方中介才会落实，城市、乡村、边地和流动人群不可能在同一时间承受相同结果。

本条使用《元史》〈本纪〉及相关志；《元朝秘史》；《元典章》或《通制条格》相关条，并标示“纪年框架可靠；该记录仅确证政权纪年连续，年内具体事项须由本年条另证”。这些材料能够钉住年序、命令、战役或局部地点，却不能直接推出人口、税收、身份和社会动机的总量。应区分诏令与实施、条约与边境生活、军报与伤亡统计，也不将官府的族群或行政称谓写成固定的社会本质。

从区域史角度看，该事件是资源、信息与权力重新连接的节点。不同地方会因市场、交通、宗教组织、家庭策略和环境条件而产生不同反应；材料的沉默限制我们对未被记录者所能作出的判断。

\eventanchor{YUAN-1280-REGNAL-YEAR}{1280年·蒙古／元以忽必烈在位第21年作为诏令、赋役与外交文书的纪年框架}账本条目记录蒙古／元在1280年的“蒙古／元以忽必烈在位第21年作为诏令、赋役与外交文书的纪年框架”。其背景是新岁颁历、年号和纪年是中枢与地方行政沟通的共同前提，后续影响为为同年政令、战事和地方材料提供日期锚点，不能单独推知具体政策执行。政权、军队或制度的变化须经由道路、文书、资源和地方中介才会落实，城市、乡村、边地和流动人群不可能在同一时间承受相同结果。

本条使用《元史》〈本纪〉及相关志；《元朝秘史》；《元典章》或《通制条格》相关条，并标示“纪年框架可靠；该记录仅确证政权纪年连续，年内具体事项须由本年条另证”。这些材料能够钉住年序、命令、战役或局部地点，却不能直接推出人口、税收、身份和社会动机的总量。应区分诏令与实施、条约与边境生活、军报与伤亡统计，也不将官府的族群或行政称谓写成固定的社会本质。

从区域史角度看，该事件是资源、信息与权力重新连接的节点。不同地方会因市场、交通、宗教组织、家庭策略和环境条件而产生不同反应；材料的沉默限制我们对未被记录者所能作出的判断。

\eventanchor{YUAN-1281-EVENT-179}{1281年·元组织第二次对日远征，舰队损失严重}账本条目记录元与日本在1281年的“元组织第二次对日远征，舰队损失严重”。其背景是前期边防、财政和统治联盟的压力构成直接背景，具体条件须按本条材料分辨，后续影响为改变后续资源配置与政治选择；实际执行范围和社会影响仍须分项核验。政权、军队或制度的变化须经由道路、文书、资源和地方中介才会落实，城市、乡村、边地和流动人群不可能在同一时间承受相同结果。

本条使用《元史》〈本纪〉及相关志；《元朝秘史》；《元典章》或《通制条格》相关条，并标示“编年主干可靠；细节、规模与动机须与对方材料、地方文书或考古材料互校”。这些材料能够钉住年序、命令、战役或局部地点，却不能直接推出人口、税收、身份和社会动机的总量。应区分诏令与实施、条约与边境生活、军报与伤亡统计，也不将官府的族群或行政称谓写成固定的社会本质。

从区域史角度看，该事件是资源、信息与权力重新连接的节点。不同地方会因市场、交通、宗教组织、家庭策略和环境条件而产生不同反应；材料的沉默限制我们对未被记录者所能作出的判断。

\eventanchor{YUAN-1281-REGNAL-YEAR}{1281年·蒙古／元以忽必烈在位第22年作为诏令、赋役与外交文书的纪年框架}账本条目记录蒙古／元在1281年的“蒙古／元以忽必烈在位第22年作为诏令、赋役与外交文书的纪年框架”。其背景是新岁颁历、年号和纪年是中枢与地方行政沟通的共同前提，后续影响为为同年政令、战事和地方材料提供日期锚点，不能单独推知具体政策执行。政权、军队或制度的变化须经由道路、文书、资源和地方中介才会落实，城市、乡村、边地和流动人群不可能在同一时间承受相同结果。

本条使用《元史》〈本纪〉及相关志；《元朝秘史》；《元典章》或《通制条格》相关条，并标示“纪年框架可靠；该记录仅确证政权纪年连续，年内具体事项须由本年条另证”。这些材料能够钉住年序、命令、战役或局部地点，却不能直接推出人口、税收、身份和社会动机的总量。应区分诏令与实施、条约与边境生活、军报与伤亡统计，也不将官府的族群或行政称谓写成固定的社会本质。

从区域史角度看，该事件是资源、信息与权力重新连接的节点。不同地方会因市场、交通、宗教组织、家庭策略和环境条件而产生不同反应；材料的沉默限制我们对未被记录者所能作出的判断。

\eventanchor{YUAN-1282-REGNAL-YEAR}{1282年·蒙古／元以忽必烈在位第23年作为诏令、赋役与外交文书的纪年框架}账本条目记录蒙古／元在1282年的“蒙古／元以忽必烈在位第23年作为诏令、赋役与外交文书的纪年框架”。其背景是新岁颁历、年号和纪年是中枢与地方行政沟通的共同前提，后续影响为为同年政令、战事和地方材料提供日期锚点，不能单独推知具体政策执行。政权、军队或制度的变化须经由道路、文书、资源和地方中介才会落实，城市、乡村、边地和流动人群不可能在同一时间承受相同结果。

本条使用《元史》〈本纪〉及相关志；《元朝秘史》；《元典章》或《通制条格》相关条，并标示“纪年框架可靠；该记录仅确证政权纪年连续，年内具体事项须由本年条另证”。这些材料能够钉住年序、命令、战役或局部地点，却不能直接推出人口、税收、身份和社会动机的总量。应区分诏令与实施、条约与边境生活、军报与伤亡统计，也不将官府的族群或行政称谓写成固定的社会本质。

从区域史角度看，该事件是资源、信息与权力重新连接的节点。不同地方会因市场、交通、宗教组织、家庭策略和环境条件而产生不同反应；材料的沉默限制我们对未被记录者所能作出的判断。

\eventanchor{YUAN-1283-REGNAL-YEAR}{1283年·蒙古／元以忽必烈在位第24年作为诏令、赋役与外交文书的纪年框架}账本条目记录蒙古／元在1283年的“蒙古／元以忽必烈在位第24年作为诏令、赋役与外交文书的纪年框架”。其背景是新岁颁历、年号和纪年是中枢与地方行政沟通的共同前提，后续影响为为同年政令、战事和地方材料提供日期锚点，不能单独推知具体政策执行。政权、军队或制度的变化须经由道路、文书、资源和地方中介才会落实，城市、乡村、边地和流动人群不可能在同一时间承受相同结果。

本条使用《元史》〈本纪〉及相关志；《元朝秘史》；《元典章》或《通制条格》相关条，并标示“纪年框架可靠；该记录仅确证政权纪年连续，年内具体事项须由本年条另证”。这些材料能够钉住年序、命令、战役或局部地点，却不能直接推出人口、税收、身份和社会动机的总量。应区分诏令与实施、条约与边境生活、军报与伤亡统计，也不将官府的族群或行政称谓写成固定的社会本质。

从区域史角度看，该事件是资源、信息与权力重新连接的节点。不同地方会因市场、交通、宗教组织、家庭策略和环境条件而产生不同反应；材料的沉默限制我们对未被记录者所能作出的判断。

\eventanchor{YUAN-1284-REGNAL-YEAR}{1284年·蒙古／元以忽必烈在位第25年作为诏令、赋役与外交文书的纪年框架}账本条目记录蒙古／元在1284年的“蒙古／元以忽必烈在位第25年作为诏令、赋役与外交文书的纪年框架”。其背景是新岁颁历、年号和纪年是中枢与地方行政沟通的共同前提，后续影响为为同年政令、战事和地方材料提供日期锚点，不能单独推知具体政策执行。政权、军队或制度的变化须经由道路、文书、资源和地方中介才会落实，城市、乡村、边地和流动人群不可能在同一时间承受相同结果。

本条使用《元史》〈本纪〉及相关志；《元朝秘史》；《元典章》或《通制条格》相关条，并标示“纪年框架可靠；该记录仅确证政权纪年连续，年内具体事项须由本年条另证”。这些材料能够钉住年序、命令、战役或局部地点，却不能直接推出人口、税收、身份和社会动机的总量。应区分诏令与实施、条约与边境生活、军报与伤亡统计，也不将官府的族群或行政称谓写成固定的社会本质。

从区域史角度看，该事件是资源、信息与权力重新连接的节点。不同地方会因市场、交通、宗教组织、家庭策略和环境条件而产生不同反应；材料的沉默限制我们对未被记录者所能作出的判断。

\eventanchor{YUAN-1285-REGNAL-YEAR}{1285年·蒙古／元以忽必烈在位第26年作为诏令、赋役与外交文书的纪年框架}账本条目记录蒙古／元在1285年的“蒙古／元以忽必烈在位第26年作为诏令、赋役与外交文书的纪年框架”。其背景是新岁颁历、年号和纪年是中枢与地方行政沟通的共同前提，后续影响为为同年政令、战事和地方材料提供日期锚点，不能单独推知具体政策执行。政权、军队或制度的变化须经由道路、文书、资源和地方中介才会落实，城市、乡村、边地和流动人群不可能在同一时间承受相同结果。

本条使用《元史》〈本纪〉及相关志；《元朝秘史》；《元典章》或《通制条格》相关条，并标示“纪年框架可靠；该记录仅确证政权纪年连续，年内具体事项须由本年条另证”。这些材料能够钉住年序、命令、战役或局部地点，却不能直接推出人口、税收、身份和社会动机的总量。应区分诏令与实施、条约与边境生活、军报与伤亡统计，也不将官府的族群或行政称谓写成固定的社会本质。

从区域史角度看，该事件是资源、信息与权力重新连接的节点。不同地方会因市场、交通、宗教组织、家庭策略和环境条件而产生不同反应；材料的沉默限制我们对未被记录者所能作出的判断。

\eventanchor{YUAN-1286-REGNAL-YEAR}{1286年·蒙古／元以忽必烈在位第27年作为诏令、赋役与外交文书的纪年框架}账本条目记录蒙古／元在1286年的“蒙古／元以忽必烈在位第27年作为诏令、赋役与外交文书的纪年框架”。其背景是新岁颁历、年号和纪年是中枢与地方行政沟通的共同前提，后续影响为为同年政令、战事和地方材料提供日期锚点，不能单独推知具体政策执行。政权、军队或制度的变化须经由道路、文书、资源和地方中介才会落实，城市、乡村、边地和流动人群不可能在同一时间承受相同结果。

本条使用《元史》〈本纪〉及相关志；《元朝秘史》；《元典章》或《通制条格》相关条，并标示“纪年框架可靠；该记录仅确证政权纪年连续，年内具体事项须由本年条另证”。这些材料能够钉住年序、命令、战役或局部地点，却不能直接推出人口、税收、身份和社会动机的总量。应区分诏令与实施、条约与边境生活、军报与伤亡统计，也不将官府的族群或行政称谓写成固定的社会本质。

从区域史角度看，该事件是资源、信息与权力重新连接的节点。不同地方会因市场、交通、宗教组织、家庭策略和环境条件而产生不同反应；材料的沉默限制我们对未被记录者所能作出的判断。

\subsection{制度、区域与材料边界}

\eventanchor{YUAN-1287-EVENT-180}{1287年·乃颜反叛，元廷在东北平定诸王叛乱}账本条目记录元在1287年的“乃颜反叛，元廷在东北平定诸王叛乱”。其背景是前期边防、财政和统治联盟的压力构成直接背景，具体条件须按本条材料分辨，后续影响为改变后续资源配置与政治选择；实际执行范围和社会影响仍须分项核验。政权、军队或制度的变化须经由道路、文书、资源和地方中介才会落实，城市、乡村、边地和流动人群不可能在同一时间承受相同结果。

本条使用《元史》〈本纪〉及相关志；《元朝秘史》；《元典章》或《通制条格》相关条，并标示“编年主干可靠；细节、规模与动机须与对方材料、地方文书或考古材料互校”。这些材料能够钉住年序、命令、战役或局部地点，却不能直接推出人口、税收、身份和社会动机的总量。应区分诏令与实施、条约与边境生活、军报与伤亡统计，也不将官府的族群或行政称谓写成固定的社会本质。

从区域史角度看，该事件是资源、信息与权力重新连接的节点。不同地方会因市场、交通、宗教组织、家庭策略和环境条件而产生不同反应；材料的沉默限制我们对未被记录者所能作出的判断。

\eventanchor{YUAN-1287-REGNAL-YEAR}{1287年·蒙古／元以忽必烈在位第28年作为诏令、赋役与外交文书的纪年框架}账本条目记录蒙古／元在1287年的“蒙古／元以忽必烈在位第28年作为诏令、赋役与外交文书的纪年框架”。其背景是新岁颁历、年号和纪年是中枢与地方行政沟通的共同前提，后续影响为为同年政令、战事和地方材料提供日期锚点，不能单独推知具体政策执行。政权、军队或制度的变化须经由道路、文书、资源和地方中介才会落实，城市、乡村、边地和流动人群不可能在同一时间承受相同结果。

本条使用《元史》〈本纪〉及相关志；《元朝秘史》；《元典章》或《通制条格》相关条，并标示“纪年框架可靠；该记录仅确证政权纪年连续，年内具体事项须由本年条另证”。这些材料能够钉住年序、命令、战役或局部地点，却不能直接推出人口、税收、身份和社会动机的总量。应区分诏令与实施、条约与边境生活、军报与伤亡统计，也不将官府的族群或行政称谓写成固定的社会本质。

从区域史角度看，该事件是资源、信息与权力重新连接的节点。不同地方会因市场、交通、宗教组织、家庭策略和环境条件而产生不同反应；材料的沉默限制我们对未被记录者所能作出的判断。

\eventanchor{YUAN-1288-REGNAL-YEAR}{1288年·蒙古／元以忽必烈在位第29年作为诏令、赋役与外交文书的纪年框架}账本条目记录蒙古／元在1288年的“蒙古／元以忽必烈在位第29年作为诏令、赋役与外交文书的纪年框架”。其背景是新岁颁历、年号和纪年是中枢与地方行政沟通的共同前提，后续影响为为同年政令、战事和地方材料提供日期锚点，不能单独推知具体政策执行。政权、军队或制度的变化须经由道路、文书、资源和地方中介才会落实，城市、乡村、边地和流动人群不可能在同一时间承受相同结果。

本条使用《元史》〈本纪〉及相关志；《元朝秘史》；《元典章》或《通制条格》相关条，并标示“纪年框架可靠；该记录仅确证政权纪年连续，年内具体事项须由本年条另证”。这些材料能够钉住年序、命令、战役或局部地点，却不能直接推出人口、税收、身份和社会动机的总量。应区分诏令与实施、条约与边境生活、军报与伤亡统计，也不将官府的族群或行政称谓写成固定的社会本质。

从区域史角度看，该事件是资源、信息与权力重新连接的节点。不同地方会因市场、交通、宗教组织、家庭策略和环境条件而产生不同反应；材料的沉默限制我们对未被记录者所能作出的判断。

\eventanchor{YUAN-1289-REGNAL-YEAR}{1289年·蒙古／元以忽必烈在位第30年作为诏令、赋役与外交文书的纪年框架}账本条目记录蒙古／元在1289年的“蒙古／元以忽必烈在位第30年作为诏令、赋役与外交文书的纪年框架”。其背景是新岁颁历、年号和纪年是中枢与地方行政沟通的共同前提，后续影响为为同年政令、战事和地方材料提供日期锚点，不能单独推知具体政策执行。政权、军队或制度的变化须经由道路、文书、资源和地方中介才会落实，城市、乡村、边地和流动人群不可能在同一时间承受相同结果。

本条使用《元史》〈本纪〉及相关志；《元朝秘史》；《元典章》或《通制条格》相关条，并标示“纪年框架可靠；该记录仅确证政权纪年连续，年内具体事项须由本年条另证”。这些材料能够钉住年序、命令、战役或局部地点，却不能直接推出人口、税收、身份和社会动机的总量。应区分诏令与实施、条约与边境生活、军报与伤亡统计，也不将官府的族群或行政称谓写成固定的社会本质。

从区域史角度看，该事件是资源、信息与权力重新连接的节点。不同地方会因市场、交通、宗教组织、家庭策略和环境条件而产生不同反应；材料的沉默限制我们对未被记录者所能作出的判断。

\eventanchor{YUAN-1290-REGNAL-YEAR}{1290年·蒙古／元以忽必烈在位第31年作为诏令、赋役与外交文书的纪年框架}账本条目记录蒙古／元在1290年的“蒙古／元以忽必烈在位第31年作为诏令、赋役与外交文书的纪年框架”。其背景是新岁颁历、年号和纪年是中枢与地方行政沟通的共同前提，后续影响为为同年政令、战事和地方材料提供日期锚点，不能单独推知具体政策执行。政权、军队或制度的变化须经由道路、文书、资源和地方中介才会落实，城市、乡村、边地和流动人群不可能在同一时间承受相同结果。

本条使用《元史》〈本纪〉及相关志；《元朝秘史》；《元典章》或《通制条格》相关条，并标示“纪年框架可靠；该记录仅确证政权纪年连续，年内具体事项须由本年条另证”。这些材料能够钉住年序、命令、战役或局部地点，却不能直接推出人口、税收、身份和社会动机的总量。应区分诏令与实施、条约与边境生活、军报与伤亡统计，也不将官府的族群或行政称谓写成固定的社会本质。

从区域史角度看，该事件是资源、信息与权力重新连接的节点。不同地方会因市场、交通、宗教组织、家庭策略和环境条件而产生不同反应；材料的沉默限制我们对未被记录者所能作出的判断。

\eventanchor{YUAN-1291-REGNAL-YEAR}{1291年·蒙古／元以忽必烈在位第32年作为诏令、赋役与外交文书的纪年框架}账本条目记录蒙古／元在1291年的“蒙古／元以忽必烈在位第32年作为诏令、赋役与外交文书的纪年框架”。其背景是新岁颁历、年号和纪年是中枢与地方行政沟通的共同前提，后续影响为为同年政令、战事和地方材料提供日期锚点，不能单独推知具体政策执行。政权、军队或制度的变化须经由道路、文书、资源和地方中介才会落实，城市、乡村、边地和流动人群不可能在同一时间承受相同结果。

本条使用《元史》〈本纪〉及相关志；《元朝秘史》；《元典章》或《通制条格》相关条，并标示“纪年框架可靠；该记录仅确证政权纪年连续，年内具体事项须由本年条另证”。这些材料能够钉住年序、命令、战役或局部地点，却不能直接推出人口、税收、身份和社会动机的总量。应区分诏令与实施、条约与边境生活、军报与伤亡统计，也不将官府的族群或行政称谓写成固定的社会本质。

从区域史角度看，该事件是资源、信息与权力重新连接的节点。不同地方会因市场、交通、宗教组织、家庭策略和环境条件而产生不同反应；材料的沉默限制我们对未被记录者所能作出的判断。

\eventanchor{YUAN-1292-REGNAL-YEAR}{1292年·蒙古／元以忽必烈在位第33年作为诏令、赋役与外交文书的纪年框架}账本条目记录蒙古／元在1292年的“蒙古／元以忽必烈在位第33年作为诏令、赋役与外交文书的纪年框架”。其背景是新岁颁历、年号和纪年是中枢与地方行政沟通的共同前提，后续影响为为同年政令、战事和地方材料提供日期锚点，不能单独推知具体政策执行。政权、军队或制度的变化须经由道路、文书、资源和地方中介才会落实，城市、乡村、边地和流动人群不可能在同一时间承受相同结果。

本条使用《元史》〈本纪〉及相关志；《元朝秘史》；《元典章》或《通制条格》相关条，并标示“纪年框架可靠；该记录仅确证政权纪年连续，年内具体事项须由本年条另证”。这些材料能够钉住年序、命令、战役或局部地点，却不能直接推出人口、税收、身份和社会动机的总量。应区分诏令与实施、条约与边境生活、军报与伤亡统计，也不将官府的族群或行政称谓写成固定的社会本质。

从区域史角度看，该事件是资源、信息与权力重新连接的节点。不同地方会因市场、交通、宗教组织、家庭策略和环境条件而产生不同反应；材料的沉默限制我们对未被记录者所能作出的判断。

\eventanchor{YUAN-1293-REGNAL-YEAR}{1293年·蒙古／元以忽必烈在位第34年作为诏令、赋役与外交文书的纪年框架}账本条目记录蒙古／元在1293年的“蒙古／元以忽必烈在位第34年作为诏令、赋役与外交文书的纪年框架”。其背景是新岁颁历、年号和纪年是中枢与地方行政沟通的共同前提，后续影响为为同年政令、战事和地方材料提供日期锚点，不能单独推知具体政策执行。政权、军队或制度的变化须经由道路、文书、资源和地方中介才会落实，城市、乡村、边地和流动人群不可能在同一时间承受相同结果。

本条使用《元史》〈本纪〉及相关志；《元朝秘史》；《元典章》或《通制条格》相关条，并标示“纪年框架可靠；该记录仅确证政权纪年连续，年内具体事项须由本年条另证”。这些材料能够钉住年序、命令、战役或局部地点，却不能直接推出人口、税收、身份和社会动机的总量。应区分诏令与实施、条约与边境生活、军报与伤亡统计，也不将官府的族群或行政称谓写成固定的社会本质。

从区域史角度看，该事件是资源、信息与权力重新连接的节点。不同地方会因市场、交通、宗教组织、家庭策略和环境条件而产生不同反应；材料的沉默限制我们对未被记录者所能作出的判断。

\eventanchor{YUAN-1294-EVENT-181}{1294年·忽必烈去世，铁穆耳即位}账本条目记录元在1294年的“忽必烈去世，铁穆耳即位”。其背景是前期边防、财政和统治联盟的压力构成直接背景，具体条件须按本条材料分辨，后续影响为改变后续资源配置与政治选择；实际执行范围和社会影响仍须分项核验。政权、军队或制度的变化须经由道路、文书、资源和地方中介才会落实，城市、乡村、边地和流动人群不可能在同一时间承受相同结果。

本条使用《元史》〈本纪〉及相关志；《元朝秘史》；《元典章》或《通制条格》相关条，并标示“编年主干可靠；细节、规模与动机须与对方材料、地方文书或考古材料互校”。这些材料能够钉住年序、命令、战役或局部地点，却不能直接推出人口、税收、身份和社会动机的总量。应区分诏令与实施、条约与边境生活、军报与伤亡统计，也不将官府的族群或行政称谓写成固定的社会本质。

从区域史角度看，该事件是资源、信息与权力重新连接的节点。不同地方会因市场、交通、宗教组织、家庭策略和环境条件而产生不同反应；材料的沉默限制我们对未被记录者所能作出的判断。

\eventanchor{YUAN-1294-REGNAL-YEAR}{1294年·蒙古／元以忽必烈在位第35年作为诏令、赋役与外交文书的纪年框架}账本条目记录蒙古／元在1294年的“蒙古／元以忽必烈在位第35年作为诏令、赋役与外交文书的纪年框架”。其背景是新岁颁历、年号和纪年是中枢与地方行政沟通的共同前提，后续影响为为同年政令、战事和地方材料提供日期锚点，不能单独推知具体政策执行。政权、军队或制度的变化须经由道路、文书、资源和地方中介才会落实，城市、乡村、边地和流动人群不可能在同一时间承受相同结果。

本条使用《元史》〈本纪〉及相关志；《元朝秘史》；《元典章》或《通制条格》相关条，并标示“纪年框架可靠；该记录仅确证政权纪年连续，年内具体事项须由本年条另证”。这些材料能够钉住年序、命令、战役或局部地点，却不能直接推出人口、税收、身份和社会动机的总量。应区分诏令与实施、条约与边境生活、军报与伤亡统计，也不将官府的族群或行政称谓写成固定的社会本质。

从区域史角度看，该事件是资源、信息与权力重新连接的节点。不同地方会因市场、交通、宗教组织、家庭策略和环境条件而产生不同反应；材料的沉默限制我们对未被记录者所能作出的判断。

\eventanchor{YUAN-1295-REGNAL-YEAR}{1295年·元以元成宗在位第1年作为诏令、赋役与外交文书的纪年框架}账本条目记录元在1295年的“元以元成宗在位第1年作为诏令、赋役与外交文书的纪年框架”。其背景是新岁颁历、年号和纪年是中枢与地方行政沟通的共同前提，后续影响为为同年政令、战事和地方材料提供日期锚点，不能单独推知具体政策执行。政权、军队或制度的变化须经由道路、文书、资源和地方中介才会落实，城市、乡村、边地和流动人群不可能在同一时间承受相同结果。

本条使用《元史》〈本纪〉及相关志；《元朝秘史》；《元典章》或《通制条格》相关条，并标示“纪年框架可靠；该记录仅确证政权纪年连续，年内具体事项须由本年条另证”。这些材料能够钉住年序、命令、战役或局部地点，却不能直接推出人口、税收、身份和社会动机的总量。应区分诏令与实施、条约与边境生活、军报与伤亡统计，也不将官府的族群或行政称谓写成固定的社会本质。

从区域史角度看，该事件是资源、信息与权力重新连接的节点。不同地方会因市场、交通、宗教组织、家庭策略和环境条件而产生不同反应；材料的沉默限制我们对未被记录者所能作出的判断。

\eventanchor{YUAN-1296-REGNAL-YEAR}{1296年·元以元成宗在位第2年作为诏令、赋役与外交文书的纪年框架}账本条目记录元在1296年的“元以元成宗在位第2年作为诏令、赋役与外交文书的纪年框架”。其背景是新岁颁历、年号和纪年是中枢与地方行政沟通的共同前提，后续影响为为同年政令、战事和地方材料提供日期锚点，不能单独推知具体政策执行。政权、军队或制度的变化须经由道路、文书、资源和地方中介才会落实，城市、乡村、边地和流动人群不可能在同一时间承受相同结果。

本条使用《元史》〈本纪〉及相关志；《元朝秘史》；《元典章》或《通制条格》相关条，并标示“纪年框架可靠；该记录仅确证政权纪年连续，年内具体事项须由本年条另证”。这些材料能够钉住年序、命令、战役或局部地点，却不能直接推出人口、税收、身份和社会动机的总量。应区分诏令与实施、条约与边境生活、军报与伤亡统计，也不将官府的族群或行政称谓写成固定的社会本质。

从区域史角度看，该事件是资源、信息与权力重新连接的节点。不同地方会因市场、交通、宗教组织、家庭策略和环境条件而产生不同反应；材料的沉默限制我们对未被记录者所能作出的判断。

\eventanchor{YUAN-1297-REGNAL-YEAR}{1297年·元以元成宗在位第3年作为诏令、赋役与外交文书的纪年框架}账本条目记录元在1297年的“元以元成宗在位第3年作为诏令、赋役与外交文书的纪年框架”。其背景是新岁颁历、年号和纪年是中枢与地方行政沟通的共同前提，后续影响为为同年政令、战事和地方材料提供日期锚点，不能单独推知具体政策执行。政权、军队或制度的变化须经由道路、文书、资源和地方中介才会落实，城市、乡村、边地和流动人群不可能在同一时间承受相同结果。

本条使用《元史》〈本纪〉及相关志；《元朝秘史》；《元典章》或《通制条格》相关条，并标示“纪年框架可靠；该记录仅确证政权纪年连续，年内具体事项须由本年条另证”。这些材料能够钉住年序、命令、战役或局部地点，却不能直接推出人口、税收、身份和社会动机的总量。应区分诏令与实施、条约与边境生活、军报与伤亡统计，也不将官府的族群或行政称谓写成固定的社会本质。

从区域史角度看，该事件是资源、信息与权力重新连接的节点。不同地方会因市场、交通、宗教组织、家庭策略和环境条件而产生不同反应；材料的沉默限制我们对未被记录者所能作出的判断。

\eventanchor{YUAN-1298-REGNAL-YEAR}{1298年·元以元成宗在位第4年作为诏令、赋役与外交文书的纪年框架}账本条目记录元在1298年的“元以元成宗在位第4年作为诏令、赋役与外交文书的纪年框架”。其背景是新岁颁历、年号和纪年是中枢与地方行政沟通的共同前提，后续影响为为同年政令、战事和地方材料提供日期锚点，不能单独推知具体政策执行。政权、军队或制度的变化须经由道路、文书、资源和地方中介才会落实，城市、乡村、边地和流动人群不可能在同一时间承受相同结果。

本条使用《元史》〈本纪〉及相关志；《元朝秘史》；《元典章》或《通制条格》相关条，并标示“纪年框架可靠；该记录仅确证政权纪年连续，年内具体事项须由本年条另证”。这些材料能够钉住年序、命令、战役或局部地点，却不能直接推出人口、税收、身份和社会动机的总量。应区分诏令与实施、条约与边境生活、军报与伤亡统计，也不将官府的族群或行政称谓写成固定的社会本质。

从区域史角度看，该事件是资源、信息与权力重新连接的节点。不同地方会因市场、交通、宗教组织、家庭策略和环境条件而产生不同反应；材料的沉默限制我们对未被记录者所能作出的判断。

\eventanchor{YUAN-1299-REGNAL-YEAR}{1299年·元以元成宗在位第5年作为诏令、赋役与外交文书的纪年框架}账本条目记录元在1299年的“元以元成宗在位第5年作为诏令、赋役与外交文书的纪年框架”。其背景是新岁颁历、年号和纪年是中枢与地方行政沟通的共同前提，后续影响为为同年政令、战事和地方材料提供日期锚点，不能单独推知具体政策执行。政权、军队或制度的变化须经由道路、文书、资源和地方中介才会落实，城市、乡村、边地和流动人群不可能在同一时间承受相同结果。

本条使用《元史》〈本纪〉及相关志；《元朝秘史》；《元典章》或《通制条格》相关条，并标示“纪年框架可靠；该记录仅确证政权纪年连续，年内具体事项须由本年条另证”。这些材料能够钉住年序、命令、战役或局部地点，却不能直接推出人口、税收、身份和社会动机的总量。应区分诏令与实施、条约与边境生活、军报与伤亡统计，也不将官府的族群或行政称谓写成固定的社会本质。

从区域史角度看，该事件是资源、信息与权力重新连接的节点。不同地方会因市场、交通、宗教组织、家庭策略和环境条件而产生不同反应；材料的沉默限制我们对未被记录者所能作出的判断。

\eventanchor{YUAN-1300-REGNAL-YEAR}{1300年·元以元成宗在位第6年作为诏令、赋役与外交文书的纪年框架}账本条目记录元在1300年的“元以元成宗在位第6年作为诏令、赋役与外交文书的纪年框架”。其背景是新岁颁历、年号和纪年是中枢与地方行政沟通的共同前提，后续影响为为同年政令、战事和地方材料提供日期锚点，不能单独推知具体政策执行。政权、军队或制度的变化须经由道路、文书、资源和地方中介才会落实，城市、乡村、边地和流动人群不可能在同一时间承受相同结果。

本条使用《元史》〈本纪〉及相关志；《元朝秘史》；《元典章》或《通制条格》相关条，并标示“纪年框架可靠；该记录仅确证政权纪年连续，年内具体事项须由本年条另证”。这些材料能够钉住年序、命令、战役或局部地点，却不能直接推出人口、税收、身份和社会动机的总量。应区分诏令与实施、条约与边境生活、军报与伤亡统计，也不将官府的族群或行政称谓写成固定的社会本质。

从区域史角度看，该事件是资源、信息与权力重新连接的节点。不同地方会因市场、交通、宗教组织、家庭策略和环境条件而产生不同反应；材料的沉默限制我们对未被记录者所能作出的判断。

\eventanchor{YUAN-1301-REGNAL-YEAR}{1301年·元以元成宗在位第7年作为诏令、赋役与外交文书的纪年框架}账本条目记录元在1301年的“元以元成宗在位第7年作为诏令、赋役与外交文书的纪年框架”。其背景是新岁颁历、年号和纪年是中枢与地方行政沟通的共同前提，后续影响为为同年政令、战事和地方材料提供日期锚点，不能单独推知具体政策执行。政权、军队或制度的变化须经由道路、文书、资源和地方中介才会落实，城市、乡村、边地和流动人群不可能在同一时间承受相同结果。

本条使用《元史》〈本纪〉及相关志；《元朝秘史》；《元典章》或《通制条格》相关条，并标示“纪年框架可靠；该记录仅确证政权纪年连续，年内具体事项须由本年条另证”。这些材料能够钉住年序、命令、战役或局部地点，却不能直接推出人口、税收、身份和社会动机的总量。应区分诏令与实施、条约与边境生活、军报与伤亡统计，也不将官府的族群或行政称谓写成固定的社会本质。

从区域史角度看，该事件是资源、信息与权力重新连接的节点。不同地方会因市场、交通、宗教组织、家庭策略和环境条件而产生不同反应；材料的沉默限制我们对未被记录者所能作出的判断。

\eventanchor{YUAN-1302-REGNAL-YEAR}{1302年·元以元成宗在位第8年作为诏令、赋役与外交文书的纪年框架}账本条目记录元在1302年的“元以元成宗在位第8年作为诏令、赋役与外交文书的纪年框架”。其背景是新岁颁历、年号和纪年是中枢与地方行政沟通的共同前提，后续影响为为同年政令、战事和地方材料提供日期锚点，不能单独推知具体政策执行。政权、军队或制度的变化须经由道路、文书、资源和地方中介才会落实，城市、乡村、边地和流动人群不可能在同一时间承受相同结果。

本条使用《元史》〈本纪〉及相关志；《元朝秘史》；《元典章》或《通制条格》相关条，并标示“纪年框架可靠；该记录仅确证政权纪年连续，年内具体事项须由本年条另证”。这些材料能够钉住年序、命令、战役或局部地点，却不能直接推出人口、税收、身份和社会动机的总量。应区分诏令与实施、条约与边境生活、军报与伤亡统计，也不将官府的族群或行政称谓写成固定的社会本质。

从区域史角度看，该事件是资源、信息与权力重新连接的节点。不同地方会因市场、交通、宗教组织、家庭策略和环境条件而产生不同反应；材料的沉默限制我们对未被记录者所能作出的判断。

\eventanchor{YUAN-1303-REGNAL-YEAR}{1303年·元以元成宗在位第9年作为诏令、赋役与外交文书的纪年框架}账本条目记录元在1303年的“元以元成宗在位第9年作为诏令、赋役与外交文书的纪年框架”。其背景是新岁颁历、年号和纪年是中枢与地方行政沟通的共同前提，后续影响为为同年政令、战事和地方材料提供日期锚点，不能单独推知具体政策执行。政权、军队或制度的变化须经由道路、文书、资源和地方中介才会落实，城市、乡村、边地和流动人群不可能在同一时间承受相同结果。

本条使用《元史》〈本纪〉及相关志；《元朝秘史》；《元典章》或《通制条格》相关条，并标示“纪年框架可靠；该记录仅确证政权纪年连续，年内具体事项须由本年条另证”。这些材料能够钉住年序、命令、战役或局部地点，却不能直接推出人口、税收、身份和社会动机的总量。应区分诏令与实施、条约与边境生活、军报与伤亡统计，也不将官府的族群或行政称谓写成固定的社会本质。

从区域史角度看，该事件是资源、信息与权力重新连接的节点。不同地方会因市场、交通、宗教组织、家庭策略和环境条件而产生不同反应；材料的沉默限制我们对未被记录者所能作出的判断。

\eventanchor{YUAN-1304-REGNAL-YEAR}{1304年·元以元成宗在位第10年作为诏令、赋役与外交文书的纪年框架}账本条目记录元在1304年的“元以元成宗在位第10年作为诏令、赋役与外交文书的纪年框架”。其背景是新岁颁历、年号和纪年是中枢与地方行政沟通的共同前提，后续影响为为同年政令、战事和地方材料提供日期锚点，不能单独推知具体政策执行。政权、军队或制度的变化须经由道路、文书、资源和地方中介才会落实，城市、乡村、边地和流动人群不可能在同一时间承受相同结果。

本条使用《元史》〈本纪〉及相关志；《元朝秘史》；《元典章》或《通制条格》相关条，并标示“纪年框架可靠；该记录仅确证政权纪年连续，年内具体事项须由本年条另证”。这些材料能够钉住年序、命令、战役或局部地点，却不能直接推出人口、税收、身份和社会动机的总量。应区分诏令与实施、条约与边境生活、军报与伤亡统计，也不将官府的族群或行政称谓写成固定的社会本质。

从区域史角度看，该事件是资源、信息与权力重新连接的节点。不同地方会因市场、交通、宗教组织、家庭策略和环境条件而产生不同反应；材料的沉默限制我们对未被记录者所能作出的判断。

\eventanchor{YUAN-1305-REGNAL-YEAR}{1305年·元以元成宗在位第11年作为诏令、赋役与外交文书的纪年框架}账本条目记录元在1305年的“元以元成宗在位第11年作为诏令、赋役与外交文书的纪年框架”。其背景是新岁颁历、年号和纪年是中枢与地方行政沟通的共同前提，后续影响为为同年政令、战事和地方材料提供日期锚点，不能单独推知具体政策执行。政权、军队或制度的变化须经由道路、文书、资源和地方中介才会落实，城市、乡村、边地和流动人群不可能在同一时间承受相同结果。

本条使用《元史》〈本纪〉及相关志；《元朝秘史》；《元典章》或《通制条格》相关条，并标示“纪年框架可靠；该记录仅确证政权纪年连续，年内具体事项须由本年条另证”。这些材料能够钉住年序、命令、战役或局部地点，却不能直接推出人口、税收、身份和社会动机的总量。应区分诏令与实施、条约与边境生活、军报与伤亡统计，也不将官府的族群或行政称谓写成固定的社会本质。

从区域史角度看，该事件是资源、信息与权力重新连接的节点。不同地方会因市场、交通、宗教组织、家庭策略和环境条件而产生不同反应；材料的沉默限制我们对未被记录者所能作出的判断。

\eventanchor{YUAN-1306-REGNAL-YEAR}{1306年·元以元成宗在位第12年作为诏令、赋役与外交文书的纪年框架}账本条目记录元在1306年的“元以元成宗在位第12年作为诏令、赋役与外交文书的纪年框架”。其背景是新岁颁历、年号和纪年是中枢与地方行政沟通的共同前提，后续影响为为同年政令、战事和地方材料提供日期锚点，不能单独推知具体政策执行。政权、军队或制度的变化须经由道路、文书、资源和地方中介才会落实，城市、乡村、边地和流动人群不可能在同一时间承受相同结果。

本条使用《元史》〈本纪〉及相关志；《元朝秘史》；《元典章》或《通制条格》相关条，并标示“纪年框架可靠；该记录仅确证政权纪年连续，年内具体事项须由本年条另证”。这些材料能够钉住年序、命令、战役或局部地点，却不能直接推出人口、税收、身份和社会动机的总量。应区分诏令与实施、条约与边境生活、军报与伤亡统计，也不将官府的族群或行政称谓写成固定的社会本质。

从区域史角度看，该事件是资源、信息与权力重新连接的节点。不同地方会因市场、交通、宗教组织、家庭策略和环境条件而产生不同反应；材料的沉默限制我们对未被记录者所能作出的判断。

\eventanchor{YUAN-1307-EVENT-182}{1307年·海山即位}账本条目记录元在1307年的“海山即位”。其背景是前期边防、财政和统治联盟的压力构成直接背景，具体条件须按本条材料分辨，后续影响为改变后续资源配置与政治选择；实际执行范围和社会影响仍须分项核验。政权、军队或制度的变化须经由道路、文书、资源和地方中介才会落实，城市、乡村、边地和流动人群不可能在同一时间承受相同结果。

本条使用《元史》〈本纪〉及相关志；《元朝秘史》；《元典章》或《通制条格》相关条，并标示“编年主干可靠；细节、规模与动机须与对方材料、地方文书或考古材料互校”。这些材料能够钉住年序、命令、战役或局部地点，却不能直接推出人口、税收、身份和社会动机的总量。应区分诏令与实施、条约与边境生活、军报与伤亡统计，也不将官府的族群或行政称谓写成固定的社会本质。

从区域史角度看，该事件是资源、信息与权力重新连接的节点。不同地方会因市场、交通、宗教组织、家庭策略和环境条件而产生不同反应；材料的沉默限制我们对未被记录者所能作出的判断。

\eventanchor{YUAN-1307-REGNAL-YEAR}{1307年·元以元成宗在位第13年作为诏令、赋役与外交文书的纪年框架}账本条目记录元在1307年的“元以元成宗在位第13年作为诏令、赋役与外交文书的纪年框架”。其背景是新岁颁历、年号和纪年是中枢与地方行政沟通的共同前提，后续影响为为同年政令、战事和地方材料提供日期锚点，不能单独推知具体政策执行。政权、军队或制度的变化须经由道路、文书、资源和地方中介才会落实，城市、乡村、边地和流动人群不可能在同一时间承受相同结果。

本条使用《元史》〈本纪〉及相关志；《元朝秘史》；《元典章》或《通制条格》相关条，并标示“纪年框架可靠；该记录仅确证政权纪年连续，年内具体事项须由本年条另证”。这些材料能够钉住年序、命令、战役或局部地点，却不能直接推出人口、税收、身份和社会动机的总量。应区分诏令与实施、条约与边境生活、军报与伤亡统计，也不将官府的族群或行政称谓写成固定的社会本质。

从区域史角度看，该事件是资源、信息与权力重新连接的节点。不同地方会因市场、交通、宗教组织、家庭策略和环境条件而产生不同反应；材料的沉默限制我们对未被记录者所能作出的判断。

\eventanchor{YUAN-1308-REGNAL-YEAR}{1308年·元以元武宗在位第1年作为诏令、赋役与外交文书的纪年框架}账本条目记录元在1308年的“元以元武宗在位第1年作为诏令、赋役与外交文书的纪年框架”。其背景是新岁颁历、年号和纪年是中枢与地方行政沟通的共同前提，后续影响为为同年政令、战事和地方材料提供日期锚点，不能单独推知具体政策执行。政权、军队或制度的变化须经由道路、文书、资源和地方中介才会落实，城市、乡村、边地和流动人群不可能在同一时间承受相同结果。

本条使用《元史》〈本纪〉及相关志；《元朝秘史》；《元典章》或《通制条格》相关条，并标示“纪年框架可靠；该记录仅确证政权纪年连续，年内具体事项须由本年条另证”。这些材料能够钉住年序、命令、战役或局部地点，却不能直接推出人口、税收、身份和社会动机的总量。应区分诏令与实施、条约与边境生活、军报与伤亡统计，也不将官府的族群或行政称谓写成固定的社会本质。

从区域史角度看，该事件是资源、信息与权力重新连接的节点。不同地方会因市场、交通、宗教组织、家庭策略和环境条件而产生不同反应；材料的沉默限制我们对未被记录者所能作出的判断。
